\chapter{Contexto del Negocio}

% ============================================================
\section{Perfiles de los Stakeholders}

\begin{longtable}{|L{3cm}|L{3.5cm}|L{7cm}|}
    \hline
    \textbf{Stakeholder} & \textbf{Rol} & \textbf{Interes / Expectativa} \\
    \hline
    \endfirsthead
    \hline
    \textbf{Stakeholder} & \textbf{Rol} & \textbf{Interes / Expectativa} \\
    \hline
    \endhead
    \hline
    \endfoot

    Alumnos & Solicitantes principales & Poder realizar trámites de forma rápida, sencilla y con seguimiento claro. No tener que acudir presencialmente ni enviar correos sin respuesta. \\
    \hline
    Docentes & Solicitantes secundarios & Solicitar listas de alumnos, constancias, apoyos SNI, estimulos y otros trámites de forma digitalizada. \\
    \hline
    Departamento Escolar (Control Escolar) & Personal que atiende solicitudes & Recibir solicitudes completas y organizadas, con toda la documentación necesaria desde el inicio. Reducir la carga de atención por multiples canales. \\
    \hline
    Responsable de Programa & Personal que atiende solicitudes & Atender solicitudes específicas de su programa de forma eficiente, con visibilidad clara de las pendientes. \\
    \hline
    Administrador del sistema & Configurador y supervisor & Poder configurar nuevos tipos de solicitud, formularios y plantillas sin depender del equipo de desarrollo. Tener visibilidad total y métricas. \\
    \hline
    Equipo de desarrollo & Desarrolladores & Contar con requerimientos claros, una arquitectura de microservicios bien definida y APIs documentadas para la integración. \\
    \hline
    Dirección de la Unidad Académica & Patrocinador & Mejorar la eficiencia administrativa, tener datos para la toma de decisiones y modernizar los procesos internos. \\
    \hline
\end{longtable}

% ============================================================
\section{Clases de Usuario}

\begin{longtable}{|L{3.2cm}|L{10.5cm}|}
    \hline
    \textbf{Usuario} & \textbf{Descripción} \\
    \hline
    \endfirsthead
    \hline
    \textbf{Usuario} & \textbf{Descripción} \\
    \hline
    \endhead
    \hline
    \endfoot

    Alumno & Estudiante inscrito en la universidad. Puede crear solicitudes de tipos como constancias, kardex, convalidación, apoyos y móvilidad. Puede consultar y dar seguimiento a sus propias solicitudes, y cancelar las que esten en estado ``Creada''. Si es alumno mentor, esta exento de adjuntar comprobante de pago. \\
    \hline
    Docente & Profesor adscrito a la universidad. Puede crear solicitudes de tipos como listas de alumnos, constancias, SNI y estimulos. Tiene las mismas capacidades de consulta y seguimiento que el alumno sobre sus propias solicitudes. \\
    \hline
    Control Escolar (Depto. Escolar) & Personal del departamento escolar. Atiende solicitudes asignadas a su rol. Puede cambiar el estado de las solicitudes, agregar observaciones y generar documentos PDF cuando el tipo de solicitud tiene plantilla. \\
    \hline
    Responsable de Programa & Coordinador de un programa académico. Atiende solicitudes asignadas a su rol con las mismas capacidades que Control Escolar dentro de su ámbito. \\
    \hline
    Administrador & Usuario con privilegios totales. Gestióna el catálogo de tipos de solicitud, los formularios dinámicos, las plantillas de documentos y la lista de alumnos mentores. Tiene visibilidad sobre todas las solicitudes y acceso al dashboard de métricas y reportes. \\
    \hline
\end{longtable}

% ============================================================
\section{Diagrama de Contexto}

El siguiente diagrama muestra como el Sistema de Solicitudes interactua con los actores (usuarios) y los sistemas externos.

\begin{figure}[H]
\centering
\begin{tikzpicture}[
    node distance=2.5cm,
    actor/.style={font=\small, align=center},
    system/.style={draw, circle, minimum size=3.5cm, fill=uazblue!10, font=\small\bfseries, align=center, text width=2.8cm},
    external/.style={draw, rectangle, rounded corners, minimum width=2.5cm, minimum height=1cm, fill=orange!15, font=\small, align=center},
    arrow/.style={-{Stealth}, thick},
    darrow/.style={{Stealth}-{Stealth}, thick},
    label/.style={font=\scriptsize, midway, fill=white, inner sep=1pt}
]

% Sistema central
\node[system] (sys) {Sistema de Solicitudes};

% Actores - izquierda
\node[actor, left=3.5cm of sys, yshift=2cm] (alumno) {\includegraphics[width=0.5cm]{example-image}\\Alumno};
\node[actor, left=3.5cm of sys, yshift=-2cm] (docente) {\includegraphics[width=0.5cm]{example-image}\\Docente};

% Actores - derecha
\node[actor, right=3.5cm of sys, yshift=2cm] (ce) {\includegraphics[width=0.5cm]{example-image}\\Control\\Escolar};
\node[actor, right=3.5cm of sys, yshift=-2cm] (rp) {\includegraphics[width=0.5cm]{example-image}\\Resp.\\Programa};

% Actor - arriba
\node[actor, above=2.5cm of sys] (admin) {\includegraphics[width=0.5cm]{example-image}\\Administrador};

% Sistemas externos - abajo
\node[external, below left=2.5cm and 1.5cm of sys] (auth) {Proveedor\\Autenticación\\(JWT)};
\node[external, below right=2.5cm and 1.5cm of sys] (siga) {SIGA\\(API)};
\node[external, below=3cm of sys] (smtp) {Servicio\\SMTP};

% Flechas actores
\draw[darrow] (alumno) -- node[label, above, sloped] {Crea, consulta,} node[label, below, sloped] {cancela solicitudes} (sys);
\draw[darrow] (docente) -- node[label, above, sloped] {Crea, consulta,} node[label, below, sloped] {cancela solicitudes} (sys);
\draw[darrow] (ce) -- node[label, above, sloped] {Atiende, cambia} node[label, below, sloped] {estado, genera PDF} (sys);
\draw[darrow] (rp) -- node[label, above, sloped] {Atiende, cambia} node[label, below, sloped] {estado, genera PDF} (sys);
\draw[darrow] (admin) -- node[label, right] {Configura catálogo, formularios, plantillas, mentores} (sys);

% Flechas sistemas externos
\draw[arrow] (auth) -- node[label, above, sloped] {Token JWT} (sys);
\draw[darrow] (sys) -- node[label, above, sloped] {Consulta datos} node[label, below, sloped] {por matricula} (siga);
\draw[arrow] (sys) -- node[label, right] {Envia correos} (smtp);

\end{tikzpicture}
\caption{Diagrama de contexto del Sistema de Solicitudes.}
\label{fig:contexto}
\end{figure}

\textbf{Descripción del diagrama:} El Sistema de Solicitudes se encuentra al centro. A su alrededor interactuan cinco tipos de actores y tres sistemas externos:

\begin{itemize}
    \item \textbf{Alumno y Docente:} Crean solicitudes, consultan su estado, dan seguimiento y pueden cancelarlas.
    \item \textbf{Control Escolar y Responsable de Programa:} Atienden solicitudes, cambian estados, agregan observaciones y generan documentos PDF.
    \item \textbf{Administrador:} Configura el catálogo de solicitudes, formularios, plantillas y lista de mentores.
    \item \textbf{Proveedor de Autenticación:} Emite tokens JWT que el sistema consume para identificar usuarios.
    \item \textbf{SIGA:} Provee datos completos de alumnos y docentes via API REST.
    \item \textbf{Servicio SMTP:} Recibe instrucciones del sistema para enviar notificaciónes por correo.
\end{itemize}

% ============================================================
\section{Árbol de Características}

El siguiente árbol presenta las funcionalidades principales del sistema organizadas por módulos.

\begin{figure}[H]
\centering
\resizebox{\textwidth}{!}{%
\begin{forest}
    for tree={
        draw, rounded corners, fill=uazblue!8,
        minimum height=0.7cm, minimum width=1.5cm,
        font=\footnotesize, align=center, text width=1.8cm,
        edge={thick, -Stealth},
        l sep=1cm, s sep=0.2cm,
        anchor=north
    }
    [Sistema de Solicitudes, fill=uazblue!25, text width=2.5cm
        [Gestión de Catálogo, text width=1.8cm
            [Crear tipo, text width=1.6cm]
            [Editar tipo, text width=1.6cm]
            [Eliminar tipo, text width=1.6cm]
        ]
        [Formularios Dinámicos, text width=1.8cm
            [Configurar campos, text width=1.6cm]
            [Tipos de campo, text width=1.6cm]
            [Validaciónes, text width=1.6cm]
        ]
        [Solicitudes, text width=1.8cm
            [Crear, text width=1.6cm]
            [Consultar, text width=1.6cm]
            [Cancelar, text width=1.6cm]
            [Seguimiento, text width=1.6cm]
        ]
        [Atención, text width=1.8cm
            [Listar asignadas, text width=1.6cm]
            [Cambiar estado, text width=1.6cm]
            [Generar PDF, text width=1.6cm]
        ]
        [Plantillas, text width=1.8cm
            [Gestiónar plantillas, text width=1.6cm]
            [Variables sustituibles, text width=1.6cm]
        ]
        [Mentores, text width=1.8cm
            [Agregar manual, text width=1.6cm]
            [Carga CSV, text width=1.6cm]
            [Listar / Eliminar, text width=1.6cm]
        ]
        [Métricas, text width=1.8cm
            [Dashboard, text width=1.6cm]
            [Exportar CSV/PDF, text width=1.6cm]
        ]
    ]
\end{forest}%
}
\caption{Árbol de características del Sistema de Solicitudes.}
\label{fig:árbol}
\end{figure}

\textbf{Descripción del diagrama:} El sistema se divide en siete módulos principales:
\begin{itemize}
    \item \textbf{Gestión de Catálogo:} CRUD de tipos de solicitud.
    \item \textbf{Formularios Dinámicos:} Configuración de campos por tipo.
    \item \textbf{Solicitudes:} Creación, consulta, cancelación y seguimiento por parte del solicitante.
    \item \textbf{Atención:} Listado, cambio de estado y generación de documentos por parte del personal.
    \item \textbf{Plantillas:} Gestión de plantillas de documentos con variables.
    \item \textbf{Mentores:} Administración de la lista de alumnos mentores exentos de pago.
    \item \textbf{Métricas:} Dashboard y exportación de reportes.
\end{itemize}
